% Les squelettes des alcanes en C5 et en C6 (corrigé du défi D11)
\begin{tikzpicture}[x=1cm,y=1cm]
\node[pcLegendeG,anchor=west,text=otNavy] at (0,3.3) {\ce{C5H12} : trois squelettes};
\node[anchor=west] at (0.2,2.45) {\chemfig[atom sep=1.6em]{-[:30]-[:-30]-[:30]-[:-30]}};
\node[anchor=west] at (4.4,2.45) {\chemfig[atom sep=1.6em]{-[:30](-[:90])-[:-30]-[:30]}};
\node[anchor=west] at (8.2,2.45) {\chemfig[atom sep=1.6em]{-[:0](-[:90])(-[:-90])-[:0]}};
\node[pcLegende,anchor=north] at (1.9,1.95) {pentane};
\node[pcLegende,anchor=north] at (5.9,1.95) {2-méthylbutane};
\node[pcLegende,anchor=north] at (9.7,1.95) {2,2-diméthylpropane};
\node[pcLegendeG,anchor=west,text=otGreenDark] at (0,1.25) {\ce{C6H14} : cinq squelettes};
\node[anchor=west] at (0.2,0.4) {\chemfig[atom sep=1.6em]{-[:30]-[:-30]-[:30]-[:-30]-[:30]}};
\node[anchor=west] at (4.6,0.4) {\chemfig[atom sep=1.6em]{-[:30](-[:90])-[:-30]-[:30]-[:-30]}};
\node[anchor=west] at (8.8,0.4) {\chemfig[atom sep=1.6em]{-[:30]-[:-30](-[:-90])-[:30]-[:-30]}};
\node[pcLegende,anchor=north] at (1.95,-0.25) {hexane};
\node[pcLegende,anchor=north] at (6.35,-0.25) {2-méthylpentane};
\node[pcLegende,anchor=north] at (10.55,-0.25) {3-méthylpentane};
\node[anchor=west] at (1.4,-1.2) {\chemfig[atom sep=1.6em]{-[:0](-[:90])(-[:-90])-[:0]-[:60]}};
\node[anchor=west] at (6.6,-1.2) {\chemfig[atom sep=1.6em]{-[:30](-[:90])-[:-30](-[:-90])-[:30]}};
\node[pcLegende,anchor=north] at (3.05,-1.85) {2,2-diméthylbutane};
\node[pcLegende,anchor=north] at (8.25,-1.85) {2,3-diméthylbutane};
\end{tikzpicture}
\par\smallskip\legende{Les huit squelettes carbonés, classés par longueur de chaîne principale décroissante : aucun autre arrangement n'est possible sans recréer l'un d'eux.}
